% BEGIN REV09_PELL_CATALAN
\subsection{Dilatación de Pell, momento de Catalán y retorno orientado}
\label{corpus9:iv:pell}

El momento paramétrico \(\mathcal C_2\) anterior conserva la
variable de profundidad; escribimos
\(G_{\mathrm{Cat}}=\mathcal C_2(1)\).
Su especialización siguiente recibe
la escala del tornillo espectral de Pell
\[
 V^{-1}\mathscr S V=\lambda_P i\,\mathscr S,\qquad
 \lambda_P=2+\sqrt3,\qquad \rho=\lambda_P^{-1}=2-\sqrt3.
\]
La ley de covariancia es la hipótesis operatoria de esta
especialización; se entiende en el dominio común invariante
de \(\mathscr S\) y del transporte invertible \(V\).
La unidad \(\lambda_P\) pertenece a \(\mathbb Q(\sqrt3)\);
la unidad \(3+2\sqrt2\) de la realización de compresión
pertenece a \(\mathbb Q(\sqrt2)\) y conserva su lector propio.
En su órbita transversal escalar,
\[
 z_n=z_0(\lambda_P i)^n,\qquad
 r_n=r_0\lambda_P^n,\qquad
 \theta_n=\theta_0+\frac{\pi_{\mathrm{HMT}}n}{2}.
\]
Para \(z_0\ne0\), eliminar \(n\) da la suspensión
\[
 r(\theta)=r_0\exp\!\left(
   \frac{2\log\lambda_P}{\pi_{\mathrm{HMT}}}
   (\theta-\theta_0)\right).
\]
La rama inversa tiene multiplicador \(-i\rho\). Esta
realización de escala recibe el tornillo como antecedente
operatorio y conserva su orientación.

\begin{proposition}[Lectura espectral y evaluación en la escala contractiva]
\label{corpus9:iv:pell-evaluacion}
En \(\mathcal H=\ell^2(\mathbb N_{>0})\), sean
\[
 N\delta_n=n\delta_n,\qquad U_4\delta_n=i^n\delta_n,\qquad
 Q_4=\frac{U_4-U_4^*}{2i},\quad
 C_4=\frac{U_4+U_4^*}{2}.
\]
Para \(0<r<1\) y \(\nu\in\mathbb C\), las trazas
\[
 F_\nu(r)=\operatorname{Tr}(N^{-\nu}r^NQ_4),\qquad
 H_\nu(r)=\operatorname{Tr}(N^{-\nu}r^NC_4)
\]
convergen absolutamente y localmente uniformemente en
\((\nu,r)\). Se tiene \(F_2(r)=\mathcal C_2(r)\) y
\begin{align}
 \mathcal C_2(\rho)
   &=\frac23G_{\mathrm{Cat}}
     -\frac{\pi_{\mathrm{HMT}}}{12}\log\lambda_P,\\
 \mathcal D\mathcal C_2(\rho)&=\frac{\pi_{\mathrm{HMT}}}{12},
 &\mathcal D^2\mathcal C_2(\rho)&=\frac14 .
 \label{corpus9:iv:pell-valores}
\end{align}
\end{proposition}
\begin{proof}
La base diagonal da las series con coeficientes
\(\sin(n\pi/2)\) y \(\cos(n\pi/2)\). En un compacto de
\(\nu\) y con \(r\leq r_0<1\), sus valores absolutos quedan
dominados por \(n^M r_0^n\) para algún \(M\), serie sumable.
La selección de los índices impares en la primera traza
produce exactamente \(\mathcal C_2(r)\).

La fórmula de sustracción de la tangente da
\(\rho=\tan(\pi_{\mathrm{HMT}}/12)\), como reconocimiento
angular posterior de la unidad de Pell. Para
\(0<\theta<\pi/2\), la integración por partes del momento da
\[
 \mathcal C_2(\tan\theta)
  =\theta\log(\tan\theta)-\int_0^\theta\log(\tan x)\,dx.
\]
El término de frontera en cero es nulo. Restando las series
de Fourier de \(\log(2\sin x)\) y \(\log(2\cos x)\), e
integrando mediante regularización de Abel, resulta
\[
 -\int_0^\theta\log(\tan x)\,dx
   =\sum_{\substack{n\geq1\\2\nmid n}}
          \frac{\sin(2n\theta)}{n^2}.
\]
La regularización tiene límite en la integral porque la
singularidad logarítmica en cero es integrable; la serie
integrada está dominada por \(\sum n^{-2}\).

Para \(n\) impar, sea \(\chi_{-4}(n)=(-1)^{(n-1)/2}\).
La identidad, comprobada por las seis clases impares módulo
doce, es
\[
 2\sin(n\pi_{\mathrm{HMT}}/6)
   =\chi_{-4}(n)
     +3\,\mathbf1_{3\mid n}\chi_{-4}(n/3).
\]
Su suma con peso \(n^{-2}\) da
\(\frac12(1+3/9)G_{\mathrm{Cat}}=2G_{\mathrm{Cat}}/3\).
Tomando \(\theta=\pi_{\mathrm{HMT}}/12\) y
\(\log\rho=-\log\lambda_P\), se obtiene la primera
evaluación. Las otras se siguen de
\(\mathcal D\mathcal C_2(r)=\arctan r\),
\(\mathcal D^2\mathcal C_2(r)=r/(1+r^2)\) y
\(\rho+\rho^{-1}=4\).
\end{proof}

La componente real de la misma traza conserva
\[
 H_2(r)=\frac14\sum_{k\geq1}\frac{(-r^2)^k}{k^2},
 \qquad
 \mathcal D H_2(r)=-\frac12\log(1+r^2),\qquad
 \mathcal D^2H_2(r)=-\frac{r^2}{1+r^2}.
\]
Estas fórmulas se obtienen seleccionando los índices pares y
derivando la serie absolutamente convergente. Por tanto,
\[
 \operatorname{Tr}(N^{-2}r^NU_4)=H_2(r)+i\mathcal C_2(r).
\]
La denominación
\(\operatorname{Li}_2(z)=\sum_{n\geq1}z^n/n^2\)
reconoce posteriormente esta traza como
\(\operatorname{Li}_2(ir)\). En particular,
\[
 \operatorname{Tr}(N^{-2}\rho^NU_4)
 =\frac14\operatorname{Li}_2(-\rho^2)
  +i\left(\frac23G_{\mathrm{Cat}}
       -\frac{\pi_{\mathrm{HMT}}}{12}\log\lambda_P\right).
\]
La sustitución de \(U_4\) por \(U_4^*\) conserva la parte
real e invierte la parte imaginaria.

\begin{proposition}[Composición de profundidad y orientación]
\label{corpus9:iv:pell-composicion}
Para \(0<r,s<1\), \(a,b\in\mathbb Z/4\mathbb Z\) y \(m\geq0\), sea
\(B(r,a)=r^NU_4^a\). Entonces
\[
 B(r,a)B(s,b)=B(rs,a+b),\qquad
 B(r,1)^m=r^{mN}U_4^m.
\]
En particular, \(B(r,1)^4=r^{4N}\ne I\) para \(0<r<1\).
\end{proposition}
\begin{proof}
Los dos factores son diagonales en la misma base. En
\(\delta_n\), su producto multiplica por
\(r^ni^{an}s^ni^{bn}=(rs)^ni^{(a+b)n}\).
La fórmula de las potencias se sigue por inducción, y
\(U_4^4=I\) da el cuarto retorno. El factor
\(r^{4N}\) permanece distinto de la identidad.
\end{proof}

El lector de potencias del multiplicador inverso
\(-i\rho\) es \(B(\rho,-1)\). Su momento imaginario tiene
orientación opuesta al de \(B(\rho,1)\), mientras la
profundidad amortiguada es la misma. El transporte \(V\),
el observable \(\mathscr S\) y el lector \(U_4\) conservan
sus espacios respectivos: la composición anterior se refiere
a los multiplicadores de la órbita y a su lectura diagonal.
La fase retorna en cuatro pasos y la contracción acumulada
permanece registrada.
% END REV09_PELL_CATALAN
